% ID: SEL-D0082-OA-E04
% Tanda S082.
\textbf{a)} Hay $120-15=105$ alumnos aptos. La proporción muestral es $\hat p=105/120=0{,}875$. Las frecuencias de 105 aptos y 15 no aptos permiten usar la aproximación normal. Al $99\%$, cada cola tiene probabilidad $0{,}005$ y $z_{0{,}995}\approx2{,}57583$. El intervalo para la proporción es
\[
 0{,}875\pm2{,}57583\sqrt{\frac{0{,}875\cdot0{,}125}{120}}
 \approx\boxed{[0{,}7972;\ 0{,}9528]}.
\]
\textbf{b)} Usamos la proporción de esta muestra como estimación previa para planificar el nuevo tamaño. El error inferior al $5\%$ se interpreta como margen menor que $0{,}05$:
\[
 2{,}57583\sqrt{\frac{0{,}875\cdot0{,}125}{n}}<0{,}05
 \quad\Longrightarrow\quad
 n>\frac{2{,}57583^2\cdot0{,}875\cdot0{,}125}{0{,}05^2}
 \approx290{,}277.
\]
El menor entero que satisface esta desigualdad estricta es $\boxed{291}$.
